\chapter{Overview of the Bible texts}

\paragraph{Where the book draws its texts from.} Number of references per book
of the Bible, in biblical order:

\begin{center}\small\begin{minipage}{.94\linewidth}\raggedright
\textosporlivro
\end{minipage}\end{center}

\paragraph{Texts that recur.} A text cited in several lessons is usually
central to the book's argument, and a good candidate for a deeper study in the
original, with the interlinear and the method of the exegesis notebooks.

\textosmaiscitados[2]

\begin{nota}[How the data was produced]
The references were extracted automatically from the PDF with
\texttt{ferramentas/livro.py textos}. The extraction recognises the
\enquote{Book chapter:verse} format; check it against the printed book and
correct \texttt{textos.tsv} if you find a mistake.
\end{nota}
